\section{Marco teórico y modelo propuesto}

\subsection{Servidor CI/CD: modelo M/M/1}

El servidor recibe trabajos según un proceso Poisson de tasa $\lambda$ y los
atiende de uno en uno con tiempos de servicio exponenciales de tasa $\mu$,
disciplina FIFO y cola infinita. Estas hipótesis corresponden a un modelo
M/M/1. La condición de estabilidad es $\rho=\lambda/\mu<1$. Para
$\lambda=0{,}25$ trabajos/min y $\mu=1/3$ trabajos/min se obtiene
$\rho=0{,}75<1$, por lo que el sistema es estable y existen las fórmulas
analíticas:
\begin{align}
L   &= \frac{\rho}{1-\rho}, &
L_q &= \frac{\rho^2}{1-\rho}, \\
W   &= \frac{1}{\mu-\lambda}, &
W_q &= \frac{\rho}{\mu-\lambda}.
\end{align}
La Ley de Little establece $L=\lambda W$ y $L_q=\lambda W_q$ para el estado
estacionario. La simulación permite estimar $L$, $L_q$, $W$ y $W_q$ y comparar
los intervalos de confianza con los valores exactos.

\subsection{SOC: múltiples servidores, prioridad y buffer finito}

El SOC tiene $s=3$ analistas en paralelo y una cola de capacidad $K=8$. La
capacidad total del sistema es $s+K=11$ alertas: tres en servicio y ocho en
espera. Si una alerta llega cuando los tres analistas están ocupados y los ocho
espacios de cola están llenos, se rechaza de inmediato y se contabiliza como
pérdida. Las alertas críticas (Tipo~1) tienen prioridad absoluta sobre las
preventivas (Tipo~2), pero la prioridad es \emph{no expropiativa}: una alerta
crítica no interrumpe un servicio en curso. Al liberarse un analista se atiende
primero la crítica más antigua; si no hay críticas, se atiende la preventiva
más antigua. Dentro de cada clase se respeta FIFO.

Las llegadas son procesos Poisson independientes con tasas $\lambda_1=0{,}05$
y $\lambda_2=0{,}10$ alertas/min en el caso base. El servicio Tipo~1 es
uniforme entre 15 y 25 minutos, con media $E[S_1]=20$ min. El servicio Tipo~2
es Weibull de forma $k=1{,}5$ y escala $\beta=12$ min; su media es
\begin{equation}
E[S_2]=\beta\,\Gamma\!\left(1+\frac{1}{k}\right)
      =12\,\Gamma\!\left(1+\frac{1}{1{,}5}\right)\approx 10{,}8329\ \text{min},
\end{equation}
no 12 minutos. La carga ofrecida por analista es
\begin{equation}
\rho_{\text{ofrecida}}=\frac{\lambda_1 E[S_1]+\lambda_2 E[S_2]}{s}.
\end{equation}
En el caso base vale $0{,}6944$; en el Escenario A, con
$\lambda_2=0{,}14$, vale $0{,}8389$.

A diferencia del M/M/1, el SOC no tiene una fórmula cerrada simple: los
servicios no son exponenciales, hay dos clases con prioridad y el buffer es
finito. Por ello, las métricas se estiman por simulación. La probabilidad de
bloqueo se define como el cociente entre alertas rechazadas y alertas
llegadas, globalmente y por tipo:
\begin{equation}
\loss=\frac{\text{bloqueadas}}{\text{llegadas}},\qquad
\loss^{(i)}=\frac{\text{bloqueadas}_i}{\text{llegadas}_i}.
\end{equation}
Las métricas de tiempo y de longitud de cola se calculan únicamente sobre las
alertas admitidas, tal como pide el enunciado.
